\documentclass[10pt,a4paper]{article}
\usepackage{amsmath,amssymb,amsthm,graphicx}
\usepackage{hyperref}
\usepackage{latexsym}
\usepackage{amssymb}
\usepackage{amsthm}
\usepackage{ulem}
\usepackage{enumitem}
\setlist{itemsep=0.4pt,topsep=0pt}
\parindent 0pt
\parindent 0pt
\parskip 12pt
\newcommand{\HCF}{\mathrm{hcf}}
\newcommand{\ZZ}{\mathbb{Z}}
\newcommand{\rk}{\mathrm{rk}}
\DeclareMathOperator{\len}{length}
\DeclareMathOperator{\bl}{b}
\DeclareMathOperator{\cl}{cl}
\DeclareMathOperator{\di}{d}
\DeclareMathOperator{\Cl}{Cl}
\DeclareMathOperator{\acl}{acl}
\DeclareMathOperator{\tp}{tp}
\DeclareMathOperator{\Th}{Th}
\newcommand{\Cor}{\mathcal{C}_\mathrm{or}}
\newcommand{\Cun}{\mathcal{C}_\mathrm{un}}
\newcommand{\mons}{\mathbb{M}}
\newcommand{\strong}{\mathrm{strong}}
\newcommand{\axioms}{\Delta_{\mathrm{FSTS}}}
\newtheoremstyle{note}% hnamei
{8pt}% hSpace abovei
{14pt}% hSpace belowi
{\normalfont}% hBody fonti
{}% hIndent amounti1
{\bfseries}% hTheorem head fonti
{ }% hPunctuation after theorem headi
{.5em}% hSpace after theorem headi2
{}% hTheorem head spec (can be left empty, meaning ‘normal’)i
\theoremstyle{note}
\newtheorem{thm}{Theorem}[section]
%\newtheorem{theorem}{Theorem}
\newtheorem{lemma}[thm]{Lemma} 
\newtheorem{df}[thm]{Definition}
\newtheorem{rmk}[thm]{Remark} 
\newtheorem{prop}[thm]{Proposition}
\newtheorem{fact}[thm]{Fact}
\newtheorem{cor}[thm]{Corollary}
\newtheorem{conj}[thm]{Conjecture}
\newtheorem{ex}[thm]{Example}
\newtheorem{qu}[thm]{Question}
\newtheorem*{notation}{Notation} 
\title{Three equivalent descriptions of when a single element is elementary over a generated subsystem}
\author{Silvia Barbina, Enrique Casanovas}
\date{arxiv:2411.13723v3; Creative Commons Attribution 4.0 International}
\begin{document}
\maketitle
\noindent\textit{Editorial note.} The counted statement is the author's environment at source0.tex lines 876-883 and its proof environment opens at line 884; the pairing is by adjacency and the retained proof range is the whole environment.\par\medskip
\noindent\textit{Editorial note.} The preamble is the author's own package lines, macro block and theorem-environment declarations (source lines 4-65); the author's optioned amsmath line and the fleqn-dependent layout line are not retained.\par\medskip
\noindent\textit{Editorial note.} The source carries no label-number dump and no counter values are imposed, so the printed numbering is the wrapper's default; identity is the author label, or the pinned source position where the author gave none.\par\medskip
\noindent\textit{Editorial note.} Metadata provenance: the source record of the pinned copy and the licence evidence file are both present; the licence is CC BY 4.0 with status supported, and the card's source\_version matches the value recorded in that evidence file.\par\medskip
\noindent\textit{Editorial scope.} One original result of the pinned author manuscript, counted once: the three equivalent descriptions for one element over a generated subsystem (author label May\_26\_2024\_1).\par\medskip
\noindent\textit{Allowed standard prerequisites (editorial).}\begin{itemize}
\item Closed subsystem. A = <A> means that A is generated by itself, so A is an elementary subsystem of M.. It is the ambient hypothesis.
\item Closure of a single element. cl\_<(a) = \{a\} means that the ordering never generates a below itself, which makes a the bottom element above the subsystem.. It is the first of the three conditions.
\item Saturated elementary extension. An omega-saturated extension carries realisations of the types that are consistent over the model.. It is where the binary tree of the third condition is realised.
\end{itemize}\medskip
\begin{lemma}\label{April_11_2024_2} The binary tree of height $\alpha$ has an HF-ordering where the root  $a_\emptyset$ is the smallest element instead of being the largest.
\end{lemma}

\begin{lemma}\label{June_25_2024_1.1} Assume $<$ is an HF-ordering of $B$,  $A=\cl_<(A)\subseteq B$  and  $C\subseteq B$ is an initial segment  such that  $c<a$  whenever $c\in C$ and $a\in A$. Then we can interpolate $A$ between $C$ and the rest of the order of $B$ in such a way that we obtain an HF-ordering. In other words,   $$(C, <\restriction C) +(A, <\restriction A) + (B\smallsetminus (A\cup C), <\restriction B\smallsetminus(A\cup C))$$ is an HF-ordering of $B$. Moreover, we can change the ordering
 of $A$ by any other HF-ordering of $A$ and still have an HF-ordering of $B$.
\end{lemma}

\begin{cor}\label{tree4}
	Let $<$ be an HF-ordering of the infinite STS $M$ and assume %$(M,<)$ 
	$M$ is $\text{$\omega$-saturated}$ as an $(R, <)$-structure.  For every $a\in M$    there is a binary tree $T=\{a_s\mid s\in 2^{<\omega}\}$  in $M$ such that
	\begin{enumerate}
		\item $a_s>a$  for every $s$
		\item $a_s>a_t$  if  $s\subsetneq t$.
	\end{enumerate}
\end{cor}

\begin{prop} \label{May_26_2024_1} Assume  $M$ is an infinite unconfined STS, $A=\langle A\rangle \subseteq  M$  and  $a\in M\smallsetminus A$. The following are equivalent:
	\begin{enumerate}
		\item  $\cl_<(a)=\{a\}$ for some HF-ordering of $M$ with initial segment  $A$
		\item   $A\leq Aa\leq M$
		\item There are $M^\prime\succeq M$, an HF-ordering $<^\prime$ of $M^\prime$ and a binary tree $T=\{a_s\mid s\in 2^{<\omega}\}$ in $M^\prime$ such that $\cl_{<^\prime}(A)=A$,  $a=a_\emptyset$, $T>^\prime A$ and  $a_s>^\prime a_{s0},a_{s1}$ for every~$s\in$~$ 2^{<\omega}$.
	\end{enumerate}
	Moreover, these conditions imply $a\leq M$.
\end{prop}

\begin{proof} 1 $\Rightarrow$ (2 and $a\leq M$). Fix an HF-ordering $<$ of $M$ with initial segment $A$ and such that $\cl_<(a)=\{a\}$.  Then  we can move $a$ to the bottom (showing $a\leq M$) and we can move $a$ to the top of $A$, making it the first element of $M\smallsetminus A$ (showing $Aa\leq M$).
	
	2 $\Rightarrow$ 1. Since $Aa\leq M$ and $A\leq Aa$, there is an HF-ordering of $M$ having an initial segment $Aa$, where $a$ is the last element of the segment. Since $a\not\in A$,  $\cl_<(a)=\{a\}$.
	
	1 $\Rightarrow$ 3. Let $<$ be an HF-ordering of $M$ with initial segment $A$ such that $\cl_<(a)=\{a\}$.  Let $(M^\prime,<^\prime)\succeq (M,<)$ be $\omega$-saturated. Then $<^\prime$ is an HF-ordering of $M^\prime$, $\cl_{<^\prime}(A)=A$,  $a>^\prime A$ and  $\cl_{<^\prime}(a)=\{a\}$.   By Corollary~\ref{tree4}, choose a tree $T_0=\{a_s\mid s\in 2^{<\omega}\}$ in  $M^\prime$ such that $a_s>^\prime a$  for every $s$ and such that
	$a_s>^\prime a_t$  if  $s\subsetneq t$. Then $a\cdot a_\emptyset>^\prime a_\emptyset$, since otherwise the blocks $\{a_\emptyset,a_0,a_1\}\neq \{a_\emptyset, a, a\cdot a_\emptyset\}$ contradict that $<^\prime$ is an HF-ordering. Let $b_0=a\cdot a_\emptyset$ and $B_0= \{b_0\}\cup T_0$.  Note that $\cl_{<^\prime}(B_0)=B_0\cup\{a\}$. Now we modify the order $<^\prime$, thus obtaining a new HF-ordering $<_0$ such that $<_0\restriction B_0= <^\prime \restriction B_0$ and $ <_0\restriction (M\smallsetminus B_0)= <^\prime\restriction (M\smallsetminus B_0)$. This ordering $<_0$ is obtained by moving  $B_0$  below $a$ but above $A$, in particular by declaring that $a>_0 b_0$ and $B_0>_0 A$ and, hence,  $a>_0a_s$ for every $s$. This can be done by Lemma~\ref{June_25_2024_1.1}. We have obtained an incomplete tree with root $a$ followed by $b_0$ on the left and the complete tree $T_0$ on the right. Now $b_0$ is a leaf and we should complete the tree below $b_0$. The construction must be iterated $\omega$ times in order to obtain a full binary tree.
	
	3 $\Rightarrow$ 1. Let $M^\prime$, $<^\prime$ and $T$ be as in the assumption. Notice that $\cl_{<^\prime}(T)=T$. By Lemma~\ref{April_11_2024_2} $T$ has an HF-ordering $<^{\prime\prime}$ where $a_\emptyset$ is the smallest element. By Lemma~\ref{June_25_2024_1.1} we can assume that $<^{\prime\prime} = <^\prime \restriction T$ and $A$ is an initial segment. Let $<=<^\prime\restriction M$. Then $<$ is an HF-ordering of $M$, $A$ is an initial segment and $\cl_<(a)=\{a\}$.
\end{proof}

\end{document}
